\chapter{Discussion}
\label{cha:n}
\section{Interpretation of Findings}
This project compared Logistic Regression, Random Forest, and Neural Network models for intrusion detection on the NSL-KDD dataset using both binary and multi-class classification. While detailed performance metrics are reported only for the held-out test set, an important observation across all models is the substantial drop in performance from validation to test data, indicating limited generalization to unseen traffic patterns. Across all models, performance on the official NSL-KDD test set is notably lower than on the validation split, indicating limited generalization to unseen traffic patterns. This gap is expected, as the test set deliberately includes attack types and distributions that are under-represented or absent during training. Higher-capacity models, particularly Random Forest and Neural Networks, exhibit a larger relative drop in performance, suggesting that strong validation results do not necessarily translate into robust test-set performance.

In the binary intrusion-detection task, all models achieve high overall accuracy and favorable ROC characteristics on the test set. However, from a security perspective, false negatives, attacks that are incorrectly classified as normal traffic, represent the most critical failure mode, as they correspond to undetected intrusions. Analysis of the confusion matrices and class-specific recall shows that Logistic Regression produces the highest number of false negatives among the evaluated models, reflecting a conservative decision boundary that prioritizes correct classification of normal traffic. Random Forest reduces the number of false negatives relative to Logistic Regression but still misclassifies a non-negligible fraction of attacks as normal, despite strong ROC--AUC values. The Neural Network achieves the lowest false-negative rate and the highest recall for attack traffic on the test set, indicating superior sensitivity to malicious behavior. These findings illustrate that ROC curves alone are insufficient for evaluating intrusion-detection performance, as they summarize ranking ability rather than actual operational decisions. Metrics that directly reflect missed attacks, such as recall for the attack class and the absolute number of false negatives, provide a more meaningful assessment of security effectiveness.

The generalization challenge becomes substantially more severe in the multi-class setting, where the task shifts from distinguishing normal from malicious traffic to discriminating between highly imbalanced attack categories. This setting exposes significantly larger differences between models, particularly when evaluation emphasizes attack sensitivity rather than class frequency. When the Normal class is excluded from metric averaging, attack-only macro recall drops substantially relative to standard macro recall for all models, demonstrating that apparent performance on the multi-class task is largely driven by correct classification of dominant classes rather than reliable attack-type detection.

Logistic Regression achieves the highest attack-only macro recall among the evaluated models. However, this result reflects a tendency toward over-generalization rather than precise attack discrimination. While recall reaches approximately $0.45$ for R2L and $0.57$ for U2R, these gains are accompanied by very low precision, indicating that the model frequently assigns minority attack labels to non-rare samples. This behavior is consistent with the limitations of linear decision boundaries in a feature space where rare attacks overlap substantially with normal traffic, leading to increased false positives as the model attempts to compensate for class imbalance. In contrast, Random Forest exhibits the weakest performance on attack categories despite its strong results in the binary setting. Recall for R2L ($\approx 0.002$) and U2R ($\approx 0.030$) is near zero, indicating that these attacks are almost entirely misclassified as normal traffic. The confusion matrices confirm that Random Forest prioritizes stable separation of dominant patterns while effectively ignoring minority attack signatures. This outcome can be attributed to both extreme class imbalance and feature importance-based learning, which favors attributes for frequent classes and suppresses weak but security-critical signals associated with rare attacks.

The Neural Network provides the most balanced multi-class attack detection. Although recall for R2L remains modest, detection of U2R improves substantially ($\approx 0.70$), demonstrating that higher-capacity models can exploit subtle feature interactions that are inaccessible to linear or tree-based approaches. However, this improvement is accompanied by increased confusion among non-dominant classes, reflecting unstable decision boundaries in regions of the feature space with limited training data. This trade-off explains why macro-averaged $F_{1}$-scores remain considerably lower than weighted metrics, despite improvements in minority recall.

Overall, these findings indicate that multi-class intrusion detection on NSL-KDD is constrained less by model choice than by data characteristics. Aggregate metrics obscure critical weaknesses in attack detection, while class-sensitive measures reveal that rare attacks remain difficult to identify. The results suggest that future improvements are more likely to arise from data-centric strategies, such as hierarchical detection pipelines, targeted resampling, or anomaly-based modeling, than from increased model complexity alone.

In the binary neural network setting, the default decision threshold of $0.5$ was not assumed a priori. Instead, decision thresholds were analyzed using validation-set precision--recall characteristics, and the threshold yielding the best macro-averaged $F_{1}$-score on the validation set was selected for final evaluation. With the changed decision threshold, both models achieve a better accuracy for recalling attacks.

\section{Contextualize with Literature}

The results obtained in this study are consistent with established findings in the NSL-KDD literature, while also exposing limitations that are frequently masked by aggregate performance metrics. Prior benchmark studies consistently report very high overall accuracy for tree-based models. For instance, Abrar et al. demonstrate that Random Forest, Extra Trees, and Decision Tree classifiers achieve accuracies exceeding 99\% on the NSL-KDD dataset when evaluated primarily using global accuracy and class-frequency weighted metrics \cite{abrar2020machine}. Similar conclusions are reported by Hegde et al. and Hong et al., who identify Random Forest as the strongest performer across accuracy, precision, and ROC--AUC after preprocessing and feature selection \cite{hegde2024intrusion, hong2021machine}. Our binary classification results align closely with these findings: all evaluated models achieve high accuracy and $F_{1}$-scores, confirming that ML-based IDS can effectively separate normal from attack traffic under the NSL-KDD benchmark.

In contrast, when extending the analysis to the multi-class setting and emphasizing attack-centric metrics, our results diverge from much of the prior literature. While many earlier studies report macro-averaged metrics that include the dominant normal and DoS classes, our use of attack-only macro recall reveals substantially weaker performance on rare attack categories such as R2L and U2R. This outcome is consistent with the foundational analysis by Tavallaee et al., who highlight that NSL-KDD, despite improvements over KDD’99, remains highly imbalanced and that minority attack types are inherently difficult to model reliably \cite{tavallaee2009detailed}. In our experiments, Random Forest exhibits near-zero recall for R2L despite strong overall accuracy, whereas Logistic Regression attains higher recall for rare classes at the cost of severe over-generalization and very low precision. Neural Networks provide a more balanced compromise, improving minority-class recall while increasing confusion among non-dominant attack categories.

These differences can be attributed primarily to methodological choices. Many reference studies prioritize accuracy-driven model selection, employ aggressive feature selection, or rely on ensemble configurations that implicitly favor majority classes. In contrast, our evaluation explicitly isolates attack detection performance and assesses generalization to rare attacks using attack-focused metrics. As a result, this work complements existing research by demonstrating that high headline accuracy does not necessarily imply effective intrusion detection for infrequent but security-critical attack types. The key contribution of this project therefore lies not in exceeding reported accuracy levels, but in providing a more nuanced, attack-aware assessment of IDS performance. This finding is consistent with earlier critiques showing that aggregate accuracy can mask poor detection of rare and previously unseen attack types under class imbalance \cite{tavallaee2009detailed}.

\section{Limitations and Future Work}

Despite the strong overall performance achieved by the evaluated models, several limitations limit the conclusions of this study. First, the NSL-KDD dataset remains highly imbalanced, with rare attack categories such as R2L and U2R representing only a small fraction of the data. This imbalance leads to substantial class overlap in the feature space and severely limits the ability of all evaluated models to detect minority attacks in the multi-class setting. Second, the fixed feature representation restricts discrimination between rare attacks and normal traffic, as many security-critical signals are weak and easily overshadowed by dominant patterns. Third, evaluation is limited to a single benchmark dataset. While NSL-KDD enables controlled comparison, it does not fully capture the diversity and dynamics of modern network traffic.

Future work should therefore prioritize data-focused strategies rather than increased model complexity. Feature selection techniques aimed at reducing class overlap, such as class-conditional importance ranking or embedded selection methods, could help isolate discriminative attributes for rare attacks. Hierarchical detection architectures are another promising direction. A staged pipeline could first separate normal from attack traffic, then classify broad attack categories, and finally identify specific subtypes. This approach can reduce the difficulty of direct multi-class learning under severe class imbalance. Finally, cost-sensitive or anomaly-aware learning methods could be applied. Weighted loss functions or hybrid anomaly--classification models can explicitly penalize missed U2R and R2L attacks and improve recall on these critical classes. These approaches directly address the main limitations of this study and would support more practical intrusion detection deployment.

\section{Updates based on the poster presentation feedbacks}

Following the feedback received during the poster presentation, several refinements were integrated into the study to enhance the depth and clarity of the analysis. A significant methodological update involved the additional use of one-hot encoding for categorical features, ensuring that the models could more effectively interpret non-numerical relationships within the data. To better illustrate the foundational challenges of the dataset, an additional graph was included to explicitly visualize the class imbalance issue, highlighting the disparity between dominant traffic and minority attack types. 

The evaluation framework was also strengthened by prioritizing security-critical metrics. Specifically, the analysis now incorporates attack-only macro recall to provide a more realistic assessment of model performance across imbalanced categories, moving beyond the optimistic bias of standard averages. Furthermore, the results section was updated to place a deliberate focus on false negatives. By analyzing these undetected intrusions more closely, the study offers a more nuanced perspective on the operational risks and the sensitivity required for effective real-world intrusion detection.